\section{问题重述与总体思路}

微网同时面临光伏发电的随机性、小区负载的时变性和外部电价的跨时段差异。储能能够把低价时段或光伏富余时段的能量转移到高价或净负荷较高的时段，但其容量、充放电功率和效率形成跨时段耦合。因此，本题的关键不是简单比较某个时刻``买不买电''，而是构造一个满足供需与储能状态约束的\textbf{多时段能量调度问题}。

四问具有明显递进关系。问题一给定典型日电价、负载和光伏预测，是确定性调度；问题二引入全年实际负载与光伏，使计划与真实净负荷产生偏差，并以 5 倍价格惩罚短缺；问题三允许在 6:00、12:00、18:00 获得新光伏预报并调整计划，是典型滚动时域控制；问题四再加入实时波动电价，需要同时处理``预测什么''和``用什么信息做决策''的因果性问题。基于这一结构，本文坚持``最简单充分模型''原则：只要目标和约束能够保持线性，就优先使用可获得全局最优解、易验证、可复现的线性规划，而不引入遗传算法、粒子群等随机元启发式方法。

总体研究路线如图所示：先完成数据审计与时间对齐，再依次构造确定性调度、风险修正日前计划、日内滚动调整与因果价格模型，最后通过约束可行性、敏感性和完全信息下界进行验证。

\begin{figure}[H]
\centering
\includegraphics{../figures_reaslab/fig01_framework.pdf}
\caption{统一建模框架与四问递进关系}
\end{figure}

为避免方法堆叠，四问共享同一能量平衡与储能状态模型，仅在信息集、预测层和结算规则上逐级扩展。这样能够把性能变化归因到新增信息或新增约束，而不是归因于无关的算法复杂度。

\section{数据审计与预处理}

附件 1 含典型日 144 个 10 min 时点的电价、负载与光伏预测；附件 2 含 2025 年 365 d $\times$ 144 点的实际负载与光伏；附件 3 含每天 0:00、6:00、12:00、18:00 四次未来 24 h 整点光伏预报；附件 4 含全年 365 d $\times$ 144 点波动电价。全部核心数值变量均无缺失值、负负载、负光伏或负电价。全年负载范围为 1995.718--7978.885 kW，实际光伏范围为 0--10216.2 kW，波动电价范围为 0.0076--1.7936 元/kWh。

进一步检查发现，附件 1 中 144 点的负载、光伏与价格几乎逐点等于附件 2/4 的全年均值，其最大差分别约为 $5.0\times 10^{-5}$ kW、0.043 kW 和 $5.2\times 10^{-5}$ 元/kWh，可将附件 1 视为典型日均值剖面。附件 3 同一日期的 6:00、12:00、18:00 行日期单元格实际为空，数据处理中按``每 4 行对应一天''重建日期，不把解析层显示的共享空字符串索引当作日期。

另一个需要明确的边界问题是：原始附件的时点为 0:10,\ldots,次日 0:00，而官方结果模板以连续 10 min 区间命名。为避免人为平移一列导致整日错位，本文\textbf{按列顺序一一对应}，并在全部计算与汇总中保持同一口径。

\begin{figure}[H]
\centering
\includegraphics{../figures_reaslab/fig02_load_pv.pdf}
\caption{典型日负载与光伏功率剖面}
\end{figure}

由典型日负载与光伏剖面可见，两者存在明显日内错峰，这为储能在中午光伏富余和夜间负荷较高时段之间进行能量搬移提供了空间。
